\documentclass[11pt,a4paper]{article}
\usepackage[T1]{fontenc}
\usepackage[utf8]{inputenc}
\usepackage{lmodern}
\usepackage{microtype}
\usepackage[margin=2.5cm]{geometry}
\usepackage{amsmath,amssymb,amsthm}
\usepackage{booktabs,array}
\usepackage{xcolor}
\usepackage{listings}
\usepackage[hidelinks]{hyperref}
\usepackage{xurl}
\setlength{\emergencystretch}{2em}


\lstset{basicstyle=\ttfamily\small,columns=fullflexible,keepspaces=true,frame=single,
  rulecolor=\color{gray!50},xleftmargin=1.2em,xrightmargin=1.2em,aboveskip=0.8em,belowskip=0.8em,
  literate={∃}{{$\exists$}}1 {∀}{{$\forall$}}1 {ℝ}{{$\mathbb{R}$}}1 {ℕ}{{$\mathbb{N}$}}1
           {↔}{{$\leftrightarrow$}}1 {≤}{{$\le$}}1 {∧}{{$\wedge$}}1 {→}{{$\to$}}1}
\newcommand{\lean}[1]{\texttt{#1}}

\newtheorem{theorem}{Theorem}[section]
\theoremstyle{remark}
\newtheorem*{remark}{Remark}

\title{An improved upper bound on the $K_4$ Ramsey multiplicity constant\\
and other formally verified results from OpenMath~2026}

\author{Woohyuk Kang\textsuperscript{1,2} \quad Hyunjin Lee\textsuperscript{3} \quad
Sanghyeon Lee\textsuperscript{4} \quad Youchan Oh\textsuperscript{5}\\[6pt]
\parbox{0.8\textwidth}{\centering\small\mbox{\textsuperscript{1}HTPeo} \quad \mbox{\textsuperscript{2}Kyung Hee University} \quad
\mbox{\textsuperscript{3}University of Cambridge}\\ \mbox{\textsuperscript{4}Korea University} \quad
\mbox{\textsuperscript{5}Seoul National University}}}
\date{October 2026}

\begin{document}
\maketitle
{\renewcommand{\thefootnote}{}\footnotetext{\textit{E-mail addresses:} \texttt{massis.kang@htpeo.com} (W.~Kang),
\texttt{hl728@cam.ac.uk} (H.~Lee), \texttt{ymhlsh4065@korea.ac.kr} (S.~Lee), \texttt{thomasoh0408@snu.ac.kr} (Y.~Oh).}}

\begin{abstract}
We report the results obtained by our team during OpenMath~2026, a one-week event in which only formally verified
mathematics was counted. Our main result is a weighted two-colouring template on 1024 blocks whose monochromatic-$K_4$
density shows that the $K_4$ Ramsey multiplicity constant satisfies $c_4 \le 0.0301399120$, improving the previous
upper bound $10486266368/768^4 \approx 0.0301422734$ due to McKay. The template, the lifting argument and the
existence of the limit defining $c_4$ are checked in Lean~4 with Mathlib, using only the standard axioms. We also
prove in Lean that several infinite families of cubic graphs have star chromatic index at most~5 and that every
bridgeless cubic loopless multigraph on at most 14 vertices has star chromatic index at most~6, and we reduce the
Dvo\v{r}\'ak--Mohar--\v{S}\'amal conjecture to named open hypotheses; the conjecture itself remains open. Finally,
we formalize known results attached to sixteen Erd\H{o}s problems, together with Sch\"onberger's and Petersen's
theorems, and report results on the scored tasks of the event, three of which carry Lean certificates. Most of the
mathematics and all Lean code were produced by AI systems directed by the authors.
\end{abstract}

\section{Introduction}

OpenMath~2026 (held as the \emph{Open Problems Hack at MIT}, 27~September--3~October 2026) invited teams to produce
machine-checked mathematics on open problems in one week~\cite{handbook}. Problems were posed as scored tasks on the
AutoLab platform; a passing evaluation was required but was explicitly not regarded as a proof, and results were
counted only when formalized. This paper collects what our team obtained, states precisely what is proved and by
which system, and records what was already known.

\paragraph{Ramsey multiplicity.}
Let $\mathrm{mono}_{K_4}(m)$ be the least number of monochromatic copies of $K_4$ over all two-colourings of the edges
of $K_m$, and $c_4=\lim_{m\to\infty}\mathrm{mono}_{K_4}(m)/\binom m4$. Erd\H{o}s conjectured $c_4=1/32$; Thomason
disproved this with a construction giving $c_4<0.03050$ and later improved it, and further small improvements
followed, among them one by Even-Zohar and Linial, as described by Parczyk, Pokutta, Spiegel and
Szab\'o~\cite{ppss}. They proved $c_4 \le 4551721\cdot2^{-24}\cdot3^{-2} < 0.03015$ from the blow-up sequence of a
Cayley graph on 768 vertices, and a note at the end of their paper reports that McKay, by local search starting from
that graph, found a graph with value $10486266368/768^4 \approx 0.0301422734$. The best known lower bound,
$c_4>0.0296$, is due to Grzesik, Lee, Lidick\'y and Volec using flag algebras, as cited in~\cite{ppss}. We prove the
following.

\begin{theorem}\label{thm:ramsey}
$c_4 \le \dfrac{200137750014677779172692777477}{6640289795261907019148161833841} \approx 0.030139911990.$
\end{theorem}

\noindent The improvement over McKay's value is about $2.4\times10^{-6}$; it narrows the gap to the lower bound by
about $0.4\%$ and does not determine $c_4$. The certificate is an explicit template, and the whole deduction,
including the existence of the limit, is checked in Lean (Section~\ref{sec:ramsey}).

\paragraph{Star edge-colouring.}
A star edge-colouring is a proper edge-colouring without bichromatic paths or cycles of length four. Dvo\v{r}\'ak,
Mohar and \v{S}\'amal~\cite{dms} proved that every subcubic graph has star chromatic index at most~7 and conjectured
that $6$ colours always suffice. We do not settle the conjecture. Our formal development works with loopless
multigraphs, which include the subcubic simple graphs of the conjecture. In Section~\ref{sec:dms} we prove in Lean that the
flower snarks, the Goldberg snarks, the generalized Petersen graphs $GP(n,k)$ with $k\le15$ and the M\"obius ladders
have star chromatic index at most~5, that every bridgeless cubic loopless multigraph on at most 14 vertices has star
chromatic index at most~6, and that the multigraph form of the conjecture is equivalent to a statement about
bridgeless cubic multigraphs on at least 16 vertices.

\paragraph{Formalizations and scored tasks.}
Section~\ref{sec:m2} describes Lean proofs of known results attached to sixteen Erd\H{o}s problems~\cite{erdosproblems},
stated exactly as in the formal-conjectures repository~\cite{fc}, and of Sch\"onberger's and Petersen's theorems for
Mathlib's simple graphs. Section~\ref{sec:hills} reports our results on the scored tasks, among them 471 triangles
formed by 39 lines in the Kobon triangle problem. Section~\ref{sec:ai} describes how AI systems were used, and
Section~\ref{sec:limits} lists what is not proved. All Lean sources, build logs and certificates are available~\cite{repo}.

\paragraph{Verification conventions.}
Unless stated otherwise, every theorem in this paper is a Lean~4.33.1 declaration~\cite{lean4} over Mathlib
v4.33.1~\cite{mathlib} that compiles without errors, and its \lean{\#print axioms} output lists only \lean{propext},
\lean{Classical.choice} and \lean{Quot.sound}. No stated theorem depends on \lean{sorry}, on added axioms or on
\lean{native\_decide}.

\section{The $K_4$ Ramsey multiplicity constant}\label{sec:ramsey}

\paragraph{Templates.}
A template on $n$ blocks consists of positive integer weights $w_1,\dots,w_n$ with total $Q$ and a symmetric
red/blue colouring of the pairs of blocks, including each block with itself. Replacing block $i$ by $t\,w_i$ vertices
and colouring each edge by the colour of its pair of blocks gives a colouring of $K_{tQ}$; its number of
monochromatic $K_4$ divided by $(tQ)^4/24$ is bounded by $N/Q^4$, where $N$ is the weighted count of monochromatic
quadruples of blocks. Hence $c_4\le N/Q^4$ for every template.

\paragraph{The certificate.}
Our template has $n=1024$ blocks, integer weights with total $Q=50762939$, and
$N = 200137750014677779172692777477$, so $N/Q^4$ is the fraction in Theorem~\ref{thm:ramsey}. The Lean development
proves the bound in the following form, together with the exact value of $N/Q^4$ and the symmetry of the template.

\begin{lstlisting}
theorem ramseyMultK4_limit_lt_ref :
    ∃ L : ℝ,
      Filter.Tendsto (fun m : ℕ => (minMonoK4 m : ℝ) / (m.choose 4 : ℝ))
        Filter.atTop (nhds L) ∧
      L < 10486266368 / 768 ^ 4
\end{lstlisting}

\noindent Two companion statements bound the same limit by $N/Q^4$ and give the bound for every $m\ge4$ rather
than only in the limit (\lean{ramseyMultK4\_le\_sol} and \lean{minMonoK4\_density\_le\_sol}).

\paragraph{What is checked and what is trusted.}
The count $N$ is split into $2\times1024$ per-block sums, each evaluated by the Lean kernel (\lean{decide +kernel}) on
a packed bit-mask representation; a general lemma proves that the packed evaluation equals the literal definition for
every template. The blow-up lemma, the monotonicity argument that gives the existence of the limit, and the
comparison with McKay's value are ordinary Lean proofs. The build consists of 2048 generated modules and took about
four CPU hours. Two steps lie outside the kernel: the transcription of the template from its data file into Lean,
which is performed by one script and checked by an independent one, and the agreement of the Lean definitions with
the informal notion of $c_4$, which we discuss in the accompanying packet.

\paragraph{How the template was found.}
We started from the 768-vertex Cayley graph used in~\cite{ppss}. Simulated annealing and tabu search acted on the
orbits of vertex pairs under its automorphism group; near-twin blocks were then split until 1024 blocks were reached,
the weights were optimized continuously and rounded to integers, and the result was refined by basin hopping. Every
candidate was re-scored exactly with the organisers' evaluator. Weight polishing alone and integer rounding gave no
further gain, and a re-splitting strategy that was still improving at the end of the event did not overtake the
best template. The template scored first of thirteen entries on the corresponding task at the end of the event.

\section{Star edge-colouring of subcubic multigraphs}\label{sec:dms}

We write $\chi'_s(G)$ for the star chromatic index. The following are proved in Lean for multigraphs given by
explicit edge lists.

\begin{theorem}\label{thm:families}
$\chi'_s(G)\le5$ for every flower snark $J_n$ ($n\ge5$ odd), every Goldberg snark, every generalized Petersen graph
$GP(n,k)$ with $1\le k\le15$ and $n\ge2k+1$ other than $GP(3,1)$, every M\"obius ladder, and every member of a family
of Petersen-type inflations.
\end{theorem}

\begin{theorem}\label{thm:le14}
Every bridgeless cubic loopless multigraph on at most $14$ vertices satisfies $\chi'_s(G)\le6$.
\end{theorem}

\begin{theorem}\label{thm:equiv}
Every subcubic loopless multigraph has star chromatic index at most~$6$ if and only if every bridgeless cubic loopless multigraph on at
least $16$ vertices, together with each of its leaf-extended subgraphs, is star $6$-edge-colourable
(\lean{dms\_iff\_cubic16}).
\end{theorem}

\noindent We also proved that a list of explicitly stated hypotheses implies the conjecture; each hypothesis is open,
and the reduction carries no evidence that they hold. The agreement of the edge lists in Theorem~\ref{thm:families}
with the textbook definitions of the families is checked by a script and not proved as a graph isomorphism. These
results were produced by an automated system in which language-model agents proposed lemmas, a Lean service accepted
or rejected them, and a second model family audited the accepted statements. Our attempts on the remaining core case,
cyclically 4-edge-connected cubic graphs, did not succeed; the informal studies are archived and are not claimed.

\section{Formalizations of known results}\label{sec:m2}

For each family in Table~\ref{tab:m2}, the Lean file is the formal-conjectures~\cite{fc} file at a fixed commit with
the \lean{sorry} of the target statement replaced by a proof; the statements are textually unchanged. We searched
formal-conjectures, the public pull requests of the awards repository for these problems and several public
collections of Lean proofs for earlier formalizations. Several candidates were dropped because a proof already
existed, and ``no earlier formal proof'' should be read as relative to these searches.

\begin{table}[ht]
\centering\small
\caption{Formalized families. Numbers refer to~\cite{erdosproblems}; theorem names omit the common prefix
\lean{erdos\_}$n$\lean{.variants}.}\label{tab:m2}
\begin{tabular}{@{}l>{\raggedright\arraybackslash}p{0.62\linewidth}l@{}}
\toprule
Problem & Result formalized & Weight\\\midrule
G-PM & Sch\"onberger's and Petersen's theorems for connected cubic bridgeless simple graphs & main\\
942 & \lean{limsup} & main\\
44 & greedy Sidon construction & main\\
123 & \lean{powers\_2\_3} & main\\
918 & three variants \lean{eq\_aleph\_0} & main\\
292 & \lean{mul}, \lean{two\_mul}, \lean{prime\_pow} & main\\
395 & \lean{one}: the radius-one counterexample of Carnielli and Carolino & main\\
698 & \lean{erdos\_szekeres} & main\\
939 & \lean{seven}, \lean{eight} & main\\
477 & \lean{S\_sq} and the quadratic case with $a\mid b$ (Sekanina's theorem and its generalization) & main\\
295, 703, 748, 1136 & helper lemma; \lean{zero}; \lean{lower\_bound}; \lean{multiples\_of\_three} & minor\\
358, 619, 1148 & odd-divisor identity; diameter-three completion; \lean{weaker} & minor\\
\bottomrule
\end{tabular}
\end{table}

\section{Results on the scored tasks}\label{sec:hills}

Table~\ref{tab:hills} lists our results on the AutoLab tasks of the event, with standings as read shortly before the
deadline; entries with identical scores share a rank.

\begin{table}[ht]
\centering\small
\caption{Results on the scored tasks.}\label{tab:hills}
\begin{tabular}{@{}>{\raggedright\arraybackslash}p{0.3\linewidth}>{\raggedright\arraybackslash}p{0.33\linewidth}>{\raggedright\arraybackslash}p{0.18\linewidth}l@{}}
\toprule
Task & Result & Standing & Lean\\\midrule
$K_4$ Ramsey multiplicity & Theorem~\ref{thm:ramsey} & 1st of 13 & yes\\
Kobon triangles, 39 lines & 471 triangles & 1st of 3 & yes\\
Busy Beaver, 6 states & halts after 249{,}881 steps & tied 1st of 12 & yes\\
Grothendieck witnesses & ratio $1393/985$, 80 bits & tied 1st of 9 & yes\\
Kobon triangles, 18 lines & 93 triangles & tied 1st of 15 & no\\
$3\times3$ matrix multiplication & rank 23, support 139 & tied 2nd of 11 & no\\
Collatz modular descent & 234 rules, full coverage & 5th of 8 & no\\
\bottomrule
\end{tabular}
\end{table}

\paragraph{Kobon triangles.}
Starting from the F\"uredi--Pal\'asti arrangement~\cite{fp}, which gives 468 triangles for 39 lines, simulated
annealing on the line parameters produced 39 lines with integer coefficients that bound 471 triangles; the upper
bound for 39 lines is $\lfloor 39\cdot 37/3\rfloor=481$. The certificate \lean{Kobon.kobon\_nonoverlapping} states that
there are 39 pairwise distinct lines and 471 triples of them, each bounding a non-degenerate triangle whose open
interior meets none of the lines, the 471 triangles being pairwise disjoint. The table of known values in
OEIS~A006066~\cite{oeis} lists no value for 39 lines; we have not searched the literature further and do not claim
the value as new.

\paragraph{Busy Beaver.}
The six-state, two-symbol machine
\begin{center}\lean{0LE1LA\_1LC1RH\_1RA1RF\_0RF0RD\_1RD1LB\_1LC1RD}\end{center}
halts from the blank tape after exactly 249{,}881 steps, visits all six states and leaves 554 ones on a span of 735
cells; this is proved in core Lean without Mathlib. The value equals the best one on the task. We did not prove that
it is the largest halting time below the hidden step budget: an exhaustive search organised by the strongly connected
structure of the transition graph excluded 31 of 32 classes and was stopped before finishing the last one.

\paragraph{Grothendieck constant.}
For the matrix $\left(\begin{smallmatrix}1&1\\1&-1\end{smallmatrix}\right)$ with rational unit vectors, the
certificate shows that every admissible constant in the finite-dimensional Grothendieck inequality is at least
$1393/985$, and at least $\sqrt2$ with exact vectors; Tsirelson's bound shows that no witness on this matrix can do
better. These facts are classical.

\paragraph{Other tasks.}
The remaining results reproduce known constructions; the matrix-multiplication scheme is one of the rank-23 schemes published by
Heule, Kauers and Seidl~\cite{hks}, and searches for support below 139 were unsuccessful. A further result of a team
member is the only entry on the held-out evaluation of the Ramsey task, with a slightly weaker template.

\begin{remark}
Two of us obtained a solution of Erd\H{o}s problem~\#1038, determining the infimum and the supremum of the measure of
$\{x:|f(x)|<1\}$ over monic real polynomials with all roots in $[-1,1]$, and formalized it in Lean. While doing so we
found that a proof by the same approach had been announced and that a Lean proof was already public. Our development
is an independent implementation of a known result, and we claim no credit for it; it is archived with the other
material~\cite{repo}.
\end{remark}

\section{Use of AI systems}\label{sec:ai}

The results of Section~\ref{sec:dms} came from an automated system running on one 16-core server, in which Claude
agents (Anthropic) proposed and proved lemmas, a Lean service accepted or rejected each submitted fact, and GPT models
(OpenAI) audited the accepted statements; about 786 informal facts were recorded. The Ramsey search, all Lean
certificates, the formalizations of Section~\ref{sec:m2} and the work on the scored tasks were carried out in Claude
Code and Codex sessions directed by the authors. The recorded usage was about $8.6\times10^7$ output tokens, all under
subscription plans; the usage of two authors was not recorded. Every proof in this paper was written by an AI system
and checked by the Lean kernel and by scripts. No author has reviewed the proofs line by line; the authors' role was
to choose problems, direct the systems, check the statements against the literature and decide what to claim. The
text of this paper was drafted with the help of an AI system and edited by the authors.

\section{What is and is not proved}\label{sec:limits}

Theorem~\ref{thm:ramsey} improves the upper bound on $c_4$ numerically; $c_4$ is not determined, and we have not
checked whether an improvement on McKay's value has appeared since 2024. The Dvo\v{r}\'ak--Mohar--\v{S}\'amal
conjecture is open; Theorem~\ref{thm:equiv} and the reduction to named hypotheses do not bring it measurably closer
to a proof. The formalizations of Section~\ref{sec:m2} concern known mathematics, and the claim that they are the
first formal proofs rests on searches whose limits we have stated. The evaluators of the scored tasks are not
formalized; where we compare a certificate with an evaluation, we compare with a transcription of the evaluator or
with its signed report. Natural next steps are a search for templates on more than 1024 blocks, an extension of
Theorem~\ref{thm:le14} to larger orders, and a formal treatment of the cyclically 4-edge-connected case.

\paragraph{Acknowledgements.}
We thank the organisers of OpenMath~2026 and the AutoLab team for the problems and the evaluation infrastructure. The
Erd\H{o}s statements and definitions are those of the formal-conjectures project~\cite{fc}; our work also relies on
Lean~\cite{lean4} and Mathlib~\cite{mathlib}.

\begin{thebibliography}{99}\small
\bibitem{dms} Z.~Dvo\v{r}\'ak, B.~Mohar and R.~\v{S}\'amal, Star chromatic index, \emph{J.~Graph Theory} 72
(2013), no.~3, 313--326, doi:10.1002/jgt.21644; arXiv:1011.3376.
\bibitem{erdosproblems} T.~F.~Bloom, Erd\H{o}s problems, \url{https://www.erdosproblems.com}, accessed 4 October 2026.
\bibitem{fc} Google DeepMind, Formal Conjectures, \url{https://github.com/google-deepmind/formal-conjectures},
commit \texttt{df3f12d7}.
\bibitem{fp} Z.~F\"uredi and I.~Pal\'asti, Arrangements of lines with a large number of triangles, \emph{Proc.
Amer. Math. Soc.} 92 (1984), 561--566, doi:10.1090/S0002-9939-1984-0760946-2.
\bibitem{handbook} Open Problems Hack at MIT (OpenMath~2026), Competition and Judging Handbook,
\url{https://rsihouse.ai/openmath/handbook.pdf}.
\bibitem{hks} M.~J.~H.~Heule, M.~Kauers and M.~Seidl, New ways to multiply $3\times3$-matrices, arXiv:1905.10192,
2019.
\bibitem{lean4} L.~de~Moura and S.~Ullrich, The Lean~4 theorem prover and programming language, in \emph{Automated
Deduction -- CADE~28}, Lecture Notes in Comput. Sci., Springer, 2021, 625--635, doi:10.1007/978-3-030-79876-5\_37.
\bibitem{mathlib} The mathlib Community, The Lean mathematical library, in \emph{Proceedings of the 9th ACM SIGPLAN
International Conference on Certified Programs and Proofs (CPP~2020)}, ACM, 2020, 367--381,
doi:10.1145/3372885.3373824; arXiv:1910.09336.
\bibitem{oeis} OEIS Foundation, Kobon triangles: maximal number of nonoverlapping triangles that can be formed from
$n$ lines drawn in the plane, Sequence A006066, \url{https://oeis.org/A006066}, accessed 4 October 2026.
\bibitem{ppss} O.~Parczyk, S.~Pokutta, C.~Spiegel and T.~Szab\'o, New Ramsey multiplicity bounds and search
heuristics, \emph{Found. Comput. Math.} 25 (2025), no.~5, 1777--1814, doi:10.1007/s10208-024-09675-6; arXiv:2206.04036.
\bibitem{repo} Lean sources, certificates and records of this work,
\url{https://github.com/lavaskiller/openmath-2026-htpeo}; state at the deadline: commit \texttt{76c63b8}.
\end{thebibliography}

\end{document}
